% Estilo do curriculo PT: replica a estrutura do
% HerberthCaldeira-Curriculo-2025.pdf (cabecalho horizontal + coluna
% principal a esquerda / sidebar a direita, titulos em azul, sans-serif).
% Arquivo separado de common/style.tex para nao afetar en/resume.tex.

\usepackage[a4paper,margin=1.6cm]{geometry}
\usepackage[T1]{fontenc}
\usepackage[utf8]{inputenc}
\usepackage{lmodern}
\renewcommand{\familydefault}{\sfdefault}
\usepackage{array}
\usepackage{enumitem}
\usepackage{xcolor}
\usepackage{hyperref}
\usepackage{paracol}

\definecolor{accent}{HTML}{1F5C99}
\definecolor{textgray}{gray}{0.45}

\pagestyle{empty}
\setlength{\parindent}{0pt}
\raggedright

\hypersetup{colorlinks=true, urlcolor=black, linkcolor=black}

% Coluna principal (esquerda, mais larga) + sidebar (direita).
\columnratio{0.64}
\setlength{\columnsep}{26pt}

% Titulo de secao: negrito, maiusculas, cor de destaque, sem regra.
% Comando manual (nao usa \section/titlesec) porque \@startsection
% conflita com o controle de colunas do paracol.
\newcommand{\cvsection}[1]{%
  \par\vspace{12pt}%
  {\color{accent}\bfseries\MakeUppercase{#1}}\par\vspace{5pt}%
}

% Cabecalho: nome + subtitulo a esquerda, contato alinhado a direita.
\newcommand{\cvheader}[4]{%
  % #1 nome, #2 subtitulo (cargo, idade), #3 localizacao, #4 telefone+email (ja formatado)
  \noindent
  \begin{minipage}[t]{0.62\textwidth}
    {\LARGE\bfseries #1}\par\vspace{2pt}
    {\large #2}
  \end{minipage}%
  \hfill
  \begin{minipage}[t]{0.34\textwidth}
    \raggedleft
    #3\\
    \textbf{#4}
  \end{minipage}
  \par\vspace{18pt}
}

% Entrada de experiencia/educacao: empresa/instituicao em negrito,
% cargo/curso em italico, periodo em cinza na linha de baixo.
\newcommand{\cventry}[3]{%
  \textbf{#1} --- \textit{#2}\par
  {\small\color{textgray}#3}\par\vspace{4pt}%
}

% Listas compactas
\setlist[itemize]{leftmargin=*, itemsep=3pt, topsep=3pt, parsep=0pt}
